\section{Introduction}
\label{sec:intro}

The rapid adoption of Kubernetes as the de-facto container orchestration
platform has introduced a complex, dynamic attack surface.
Unlike traditional datacenter environments, Kubernetes clusters feature
ephemeral workloads, lateral-movement-friendly overlay networks, and
shared kernel namespaces that static signature-based intrusion detection
systems (IDS) struggle to monitor effectively~\cite{sysdig2023threat}.

Anomaly-based detection offers a complementary approach: by learning what
\textit{normal} traffic looks like, it can flag deviations without relying
on prior knowledge of attack signatures.
However, real-world Kubernetes traffic is highly heterogeneous, combining
short-lived DNS probes, long-duration TLS sessions, and bursty microservice
calls within the same namespace.
A single model class is unlikely to capture all anomalous patterns.

\textbf{Contributions.}
\begin{enumerate}
  \item We propose a \emph{weighted ensemble} of Isolation Forest (weight 0.4)
        and an LSTM Autoencoder (weight 0.6) that exploits both structural
        and temporal anomalies in network flows.
  \item We describe a fully integrated Kubernetes security stack—Istio mTLS,
        OPA Gatekeeper, Falco, and an ELK SIEM—that provides defence-in-depth
        around the ML detector.
  \item We provide an open-source implementation that runs on a single
        8-GB workstation, making the system accessible to researchers and
        small-to-medium enterprises.
  \item We evaluate the system on UNSW-NB15~\cite{moustafa2015unsw} and
        CIC-IDS2017~\cite{sharafaldin2018toward} and report precision, recall,
        F1, and AUC-ROC for each component and the ensemble.
\end{enumerate}

The remainder of the paper is structured as follows.
Section~\ref{sec:related} surveys related work.
Section~\ref{sec:design} describes the system architecture.
Section~\ref{sec:ml} details the ML pipeline.
Section~\ref{sec:eval} presents experimental results.
Section~\ref{sec:discussion} discusses limitations and future work.
Section~\ref{sec:conclusion} concludes.
